%%%%%%%%%%%%%%%%%%%%%%%%%%%%%%%%%%%%%%%%%%%%%%%%%%%%%%%%%%%%%%%%%%%%%%%%%%%%
% AGUJournalTemplate.tex: this template file is for articles formatted with LaTeX
%
% This file includes commands and instructions
% given in the order necessary to produce a final output that will
% satisfy AGU requirements, including customized APA reference formatting.
%
% You may copy this file and give it your
% article name, and enter your text.
%
% guidelines and troubleshooting are here:

%% To submit your paper:
\documentclass{agujournal2019}

\usepackage{url} %this package should fix any errors with URLs in refs.
\usepackage{lineno}
\usepackage[inline]{trackchanges} %for better track changes. finalnew option will compile document with changes incorporated.
\usepackage{soul}
\usepackage{amsmath}
\usepackage{amssymb}


\linenumbers

\draftfalse

\journalname{JGR: Earth Surface, Confidential manuscript under review}


\begin{document}

\title{Constraining the Distribution of 3D Fractal Structures in Mud Flocs}
\authors{
Brayden Noh\affil{1,2},
Justin A.~Nghiem\affil{1,3},
Kimberly Litwin Miller\affil{1},
Dongchen Wang\affil{1},
Michael P.~Lamb\affil{1}
}

\affiliation{1}{Division of Geological and Planetary Sciences, California Institute of Technology, Pasadena, CA 91125, USA}
\affiliation{2}{Department of Earth and Planetary Sciences, Harvard University, Cambridge, MA 02138, USA}
\affiliation{3}{St. Anthony Falls Laboratory, University of Minnesota, Minneapolis, MN 55414, USA}

\correspondingauthor{Brayden Noh}{braydennoh@fas.harvard.edu}


\begin{keypoints}
\item Aggregating floc data overestimates 3D fractal dimension due to structural diversity.
\item Stratifying by 2D box-counting dimension reveals trends with shear velocity and organic matter.
\item Revised fractal dimensions indicate lower settling rates than conventional aggregated estimates.
\end{keypoints}


\begin{abstract}
Mud builds coastal landscapes and carries pollutants and carbon, yet predicting its transport is difficult because mud aggregates into flocs that control its settling rate in river, estuarine, and marine settings. Flocs are typically characterized using fractal theory with a constant fractal dimension ($n_f = 2.0$), which is inferred from aggregating floc size--settling velocity data, despite a diversity of complex fractal structures. We conducted experiments on freshwater flocs under varied shear stress and organic concentration to characterize this diversity. Rather than the conventional approach of aggregating settling data to infer $n_f$, we grouped measurements by the two-dimensional (2D) box-counting dimension derived from particle tracking, thereby partially resolving structural variability. This stratified method reveals that $n_f$ increases with shear velocity and decreases with organic matter content, trends that are obscured when data are aggregated. The stratified analysis resolves $n_f$ variations of 0.2--0.35 units across experimental conditions, whereas the conventional aggregated approach reduces this range by roughly a factor of five. The inferred primary particle diameter $d_p$ co-varies with these conditions, and the concurrent variation of $n_f$, $d_p$, and floc size introduces competing effects on settling velocity. Our approach yields lower $n_f$ values (1.6--2.1) than the canonical range, which propagates nonlinearly into settling rates, implying that mud flocs are more porous and settle more slowly than expected. We also find that the 2D-to-3D $n_f$ relationship can be described by a nonlinear model, consistent with empirical conversions for synthetic aggregates. Overall, our results indicate that mud retention in floodplains and deltas is less efficient than current floc models assume, at least in the freshwater conditions we analyzed, with implications for sediment budgets, contaminant dispersal, and coastal landscape evolution.
\end{abstract}

\section*{Plain Language Summary}
Mud particles in rivers tend to clump together into clusters called flocs. How quickly these flocs sink determines where mud is deposited, shaping the growth of river deltas, floodplains, and coastlines. Predicting this settling is challenging because flocs vary widely in structure: some are dense and compact, while others are porous and loosely bound. Most studies assume that all flocs have a single, uniform structure, which oversimplifies their behavior. In our experiments, we generated flocs in freshwater under controlled conditions with different flow strengths and amounts of organic matter. Using high-resolution imaging, we measured both their structure and settling speed. We found that faster flows tend to produce smaller but more compact flocs, while organic matter encourages the formation of more open structures that settle more slowly. By grouping flocs according to their measured structural complexity, we revealed trends that are hidden when data are averaged together. This new approach reduces bias in estimating settling rates and improves predictions of how mud moves through rivers and into estuaries and oceans, where salinity may further modify aggregation beyond the freshwater trends isolated here.

\section{Introduction}

Fine-grained sediment, or mud, defined by particles smaller than 62.5~\textmu m, dominates fluvial sediment loads and represents the dominant material deposited in lowland and deltaic environments \cite{healy2002muddy, kranenburg1994fractal, lamb2020mud, nghiem2026flocculated}. The dynamics of mud transport and deposition therefore exert a first-order control on the long-term evolution of deltas, floodplains, and estuaries \cite{nghiem2022mechanistic, winterwerp2002flocculation}, as well as on water quality, contaminant fluxes, and the success of restoration interventions \cite{droppo1997freshwater, esposito2017efficient, kemp2016enhancing, liss1996floc}. Mud also plays a central role in the global carbon cycle by providing reactive mineral surfaces for organic carbon adsorption and long-term preservation \cite{bianchi2024anthropogenic, blair2012fate}.

The geomorphic and geochemical importance of mud necessitates a mechanistic framework for particle transport and settling to enable predictive modeling. Settling velocity is a key parameter in sediment-transport equations and is commonly estimated using Stokes’ law, which gives the terminal settling velocity ($w_s$) of smooth, impermeable spheres at low Reynolds numbers \cite{bird2006transport, stokes1851effect}:
\begin{equation}
w_s = \frac{g\,(\rho_p - \rho_f)\,d^2}{18\,\rho_f\,\nu}
\label{eq:stokes}
\end{equation}
where $g = 9.81\ \mathrm{m\,s^{-2}}$ is gravitational acceleration, $\rho_p\ (\mathrm{kg\,m^{-3}})$ and $\rho_f\ (\mathrm{kg\,m^{-3}})$ are the densities of the particle and fluid, $d\ (\mathrm{m})$ is particle diameter, and $\nu\ (\mathrm{m^2\,s^{-1}})$ is the kinematic viscosity of the fluid.

However, when applied to fine-grained sediment, Stokes’ law yields predictions at odds with field observations. Consider a single $1~\text{\textmu m}$ mud particle in a $5$~m-deep river flowing at $1~\mathrm{m\,s^{-1}}$. Using Stokes’ law produces a settling velocity of just $9\times10^{-8}~\mathrm{m\,s^{-1}}$, implying a deposition time greater than $600$~days and a transport distance exceeding $55{,}000$~km downstream. This hypothetical scenario conflicts sharply with field observations of mud deposition in fluvial and coastal settings \cite{lamb2020mud, mason2019differential, nicholas1996significance}.

This discrepancy indicates that mud settles far more rapidly than predicted for isolated particles, pointing to additional physical processes. The deviation arises because suspended mud particles flocculate, forming larger aggregates composed of individual primary particles that settle much faster than their individual constituents \cite{droppo1997freshwater, johnson1996settling, kranenburg1994fractal}. These flocculated aggregates, or flocs, are highly porous rather than dense, impermeable spheres, so their effective density decreases with increasing size \cite{kranenburg1994fractal, Meakin1992Aggregation}. Because flocs exhibit approximately self-similar fractal organization \cite{family2012kinetics, jullien1987aggregation}, their hierarchical structure can be described by a fractal mass--size scaling \cite{Meakin1992Aggregation}:
\begin{equation}
N \sim \left( \frac{d_f}{d_p} \right)^{n_f},
\label{eq:fractaltheory}
\end{equation}
where $N$ is the number of primary particles (dimensionless), $d_f$ is the floc diameter (m), $d_p$ is the primary particle diameter (m), and $n_f$ is the fractal dimension (dimensionless).

Building on Eq.~\eqref{eq:fractaltheory}, the fraction of the floc volume occupied by solid grains is the ratio of solid volume to the enclosing floc volume. Because the solid volume scales with $N$ times the monomer volume while the enclosure scales with $d_f^{3}$, the solid volume fraction $\phi$ decreases with size approximately as $(d_f/d_p)^{n_f-3}$. For a porous aggregate whose pores are filled by ambient water, the floc’s submerged specific gravity equals this solid fraction times the solid’s submerged specific gravity. Denoting submerged specific gravity by $R = (\rho - \rho_w)/\rho_w$, this yields:
\begin{equation}
R_f = R_s \left( \frac{d_f}{d_p} \right)^{n_f-3},
\label{eq:eq3}
\end{equation}
where $R_f$ and $R_s$ are the submerged specific gravity of the floc and primary particles, respectively (both dimensionless), $\rho$ is particle density (kg\,m$^{-3}$), and $\rho_w$ is water density (kg\,m$^{-3}$). Since a dense solid sphere has volume that scales with the cube of its diameter, its fractal dimension is $n_f = 3$, and Eq.~\eqref{eq:eq3} reduces to $R_f = R_s$. Naturally formed flocs are more porous with $1 \le n_f < 3$, yielding $R_f < R_s$. Accordingly, floc effective (excess) density decreases with floc size. For a fractal aggregate, the solids mass scales as $M_f \propto d_f^{n_f}$ whereas the envelope volume scales as $V_f \propto d_f^{3}$, implying $\Delta \rho_f \propto M_f / V_f \propto d_f^{n_f-3}$ \cite{kranenburg1994fractal}.

Substituting Eq.~\eqref{eq:eq3} into Eq.~\eqref{eq:stokes} gives an explicit expression for the settling velocity of a fractal floc \cite{strom2011explicit, winterwerp1998simple}:
\begin{equation}
w_s = \frac{g\,R_s}{b_1\,\nu\,d_p^{\,n_f-3}}\, d_f^{\,n_f-1},
\label{eq:floc_settling}
\end{equation}
where $w_s$ is the settling velocity (m\,s$^{-1}$), $g$ is gravitational acceleration (m\,s$^{-2}$), $\nu$ is kinematic viscosity (m$^{2}$\,s$^{-1}$), and $b_1$ is a dimensionless shape factor (with $b_1 = 18$ for smooth spheres).

In practice, applying Eq.~\eqref{eq:floc_settling} requires an appropriate fractal dimension $n_f$, which captures how floc structure modifies settling velocity. Experimental and field studies often assume $n_f \approx 2$ from empirical $w_s$–$d_f$ relationships, values reported primarily for saline estuarine/coastal flocs where electrochemical interactions promote more compact aggregates; freshwater flocs more often show lower $n_f$ ($\lesssim 2$), with overlap depending on turbulence and composition \cite{winterwerp1998simple, winterwerp2006heuristic}. With all other parameters held fixed, Eq.~\eqref{eq:floc_settling} gives $w_s \propto d_f^{\,n_f-1}$ \cite{johnson1996settling, li1989fractal}. Measurements of $w_s$ and $d_f$ across a floc population therefore allow $n_f$ to be inferred from the log–log slope of the $w_s$–$d_f$ relation \cite{kranenburg1994fractal, strom2011explicit, winterwerp1998simple, winterwerp2002flocculation, winterwerp2004introduction, winterwerp2006heuristic}. In general, this inference assumes that floc permeability is independent of floc size. If permeability instead varies systematically with $d_f$, the observed scaling relation may reflect both fractal structure and permeability effects \cite{strom2011explicit, nghiem2026flocculated}.

Using a constant population-wide value for the fractal dimension is convenient but restrictive. \citeA{khelifa2006models} showed that treating $n_f$ as constant fails to reproduce the observed spread in settling velocities and argued that $n_f$ can decrease with floc size, motivating models that prescribe a function $n_f(d_f)$, often as a power law, to represent within-population variability. In practice, however, most experimental and image-based field studies measure only floc size $d_f$ and settling velocity $w_s$ \cite{smith2015image, nghiem2024testing}, so the fractal dimension is not independently observed and functions for $n_f$ cannot be directly tested. Moreover, models for $n_f(d_f)$ implicitly assume a unique fractal dimension for a given floc size \cite{khelifa2006models}, whereas natural floc populations can contain multiple flocs with comparable $d_f$ but distinct internal structure, and therefore distinct $n_f$ \cite{maggi2007variable}. As a result, standard log–log regressions of $w_s$ on $d_f$ treat structurally heterogeneous populations as homogeneous, and pooling such subpopulations can misattribute between-group differences to size scaling itself, producing biased estimates of $n_f$.

To illustrate this effect, consider a population comprising three subgroups of flocs, each defined by a distinct fractal dimension ($n_f = 1.5$, $2.0$, and $2.5$) but formed from primary particles of the same diameter ($d_p = 10^{-5}$ m). According to Eq.~\eqref{eq:floc_settling}, the slope of the log–log relationship between settling velocity and floc diameter for each subgroup is $n_f - 1$. When analyzed separately, these subgroups yield distinct log–log slopes corresponding to $n_f = 1.5$, $2.0$, and $2.5$, respectively (Fig.~\ref{fig:1}). However, if these floc subgroups are combined into a single group and floc diameter ($d_f$) is correlated with the 3D fractal dimension ($n_f$), then fitting a single power-law relationship to all points can yield a substantially skewed and misleading exponent. In Fig.~\ref{fig:1}, the resulting log–log slope is $2$, implying $n_f = 3$, which corresponds to the settling behavior of smooth, impermeable spheres, not flocs. This example parallels Simpson’s Paradox \cite{simpson1951interpretation}, a statistical phenomenon in which relationships observed within individual subgroups are reversed, masked, or distorted when data are pooled. In the context of floc settling, pooling structurally distinct subpopulations can produce a global $w_s$–$d_f$ slope that implies an apparent $n_f$ inconsistent with the $n_f$ of any subgroup, so the inferred ``representative'' fractal dimension is a regression artifact rather than a physical descriptor.
\begin{figure}[htbp]
    \centering
    \includegraphics[width=0.8\textwidth]{figures/fig1.pdf}
    \caption{
    Synthetic data illustrating Simpson's Paradox in floc settling. Three floc subgroups, each with a specific 3D fractal dimension ($n_f = 1.5,\ 2.0,\ \text{and}\ 2.5$), are shown. By Eq.~\eqref{eq:floc_settling}, the slope of the log--log relationship between settling velocity ($w_s$) and floc diameter ($d_f$) is $n_f - 1$. Fitting each subgroup individually recovers slopes of 0.5, 1.0, and 1.5, as expected. In contrast, pooling all points produces a single power-law fit (red line) with a misleading slope of 2, equivalent to $n_f = 3$, characteristic of solid spheres and not flocs. All trend lines intersect at the primary particle diameter ($d_f = d_p = 10^{-5}\,\mathrm{m}$), where settling velocity is independent of fractal dimension.
    }
    \label{fig:1}
\end{figure}

To estimate fractal dimension independently of settling velocity and floc diameter, studies have employed image-based techniques that analyze in situ 2D projections of flocs \cite{jarvis2005measuring}. These approaches typically calculate a two-dimensional (2D) fractal dimension from high-resolution images, then infer the corresponding three-dimensional (3D) structure via empirical or simulation-based relationships. A common method relates the perimeter--area scaling of projected flocs to a 2D fractal dimension, using computationally generated aggregates to derive a mapping between 2D and 3D fractal dimension \cite{lee2004prediction, maggi2004method, tang2015reconstructing}. However, the relationship between 2D and 3D fractal dimension is derived from synthetic aggregates, which may not replicate the aggregation dynamics or structural variability of natural flocs. Projection-based estimates therefore risk systematic bias and, more importantly, do not recover the 3D fractal dimension that directly controls settling behavior.

Although flocculation has been widely characterized, how within-population structural heterogeneity affects settling remains uncertain, particularly under natural particle-size distributions and aggregation pathways. The central difficulty is that $w_s$–$d_f$ regressions alone do not condition on structure, so pooling distinct subpopulations can produce an apparent relation and inferred $n_f$ that is biased by data-pooling (Simpson-type) effects. To address this bias, we use images from experiments to stratify the floc population by the 2D box-counting dimension and within each subgroup, regress settling velocity $w_s$ against floc diameter $d_f$ to estimate the hydrodynamically relevant 3D fractal dimension $n_f$. This conditioning step provides a principled way to separate structural heterogeneity from size scaling without prescribing a population-level closure such as $n_f(d_f)$: that is, $n_f$ is inferred hydrodynamically within approximately homogeneous structural classes. The assumption is that flocs with similar 2D box-counting dimension have similar 3D fractal dimension, which we validate herein.

We evaluate this framework in controlled laboratory experiments designed to simulate flocs in rivers with systematic variation in shear velocity and organic matter concentration. High-resolution in situ imaging enabled quantification of floc morphology and settling trajectories. Specifically, we pursued three objectives: (1) demonstrate the existence of structurally distinct subgroups within the floc population, each characterized by a unique fractal dimension; (2) relate 2D image–derived fractal dimensions to settling-inferred 3D dimensions, and assess this relationship against theoretical predictions; and (3) estimate the effective primary particle size ($d_p$) and quantify how the hydrodynamic shape factor ($b_1$) varies as a function of $n_f$. All experiments were conducted in freshwater (salinity $\approx 0$), so the results isolate the roles of shear and organic binding and are not intended as a direct constraint for saline estuaries where salt-driven cohesion can further modify flocculation.

\section{Methodology}

\subsection{Experimental Setup}

We performed flocculation experiments in a custom mixing tank designed to provide controlled conditions (Fig.~\ref{fig:2}a). The tank comprises a cylindrical PVC pipe, 20~cm in diameter and 60~cm tall, filled with water to a depth of 40~cm. The tank was filled with fresh tap water, used without filtration or chemical modification, and experiments were conducted at room temperature. A motor-driven rotating lid, mounted above the open top and submerged just below the free surface, sets the desired shear. For each experiment, the tank was filled with water and the motor was started to establish the target shear. Particles and organic matter (described below) were then added gradually, and the tank was mixed for 30~min to allow flocs to form under sustained shear. The motor was subsequently stopped to initiate settling, and video recording began immediately and continued for the first 30~min.

Shear velocity, $u_*$, was empirically calibrated for each run by converting motor revolutions per minute to $u_*$ using a relationship derived from previous sediment-entrainment trials in \citeA{douglas2025mud}. We define $u_* = (\tau/\rho_w)^{1/2}$, where $\tau$ is a characteristic bed shear stress, and we treat $u_*$ as a characteristic measure of turbulent shear intensity because the rotating lid-driven flow field and boundary stresses are spatially heterogeneous.

During the mixing phase, the flow was fully turbulent. This is quantified by the friction Reynolds number $Re_{\tau} = u_* h / \nu$, where $h$ is the water depth and $\nu$ is the kinematic viscosity of freshwater. For $h = 0.40$~m and $\nu \approx 1.0 \times 10^{-6}\ \mathrm{m^2\,s^{-1}}$ at room temperature, the imposed shear velocities $u_* = 0.02$–$0.04$~m~s$^{-1}$ correspond to $Re_{\tau} \approx (8$–$16)\times10^{3}$, well within the fully turbulent regime for wall-bounded shear flows \cite{Pope2000, Schlichting2000}. Across the flocculation experiments, we imposed $u_* = 0.02$–$0.04$~m~s$^{-1}$ (equivalent to $\tau \approx 0.4$–$1.6$~Pa), spanning low-to-moderate shear conditions typical of suspension-dominated river channels, sufficient for sand entrainment but below those required to mobilize gravels \cite{julien2018river}.

All image-based measurements reported here were collected during the settling phase. Suspended flocs were imaged through a 3~mm-thick flow-through slit built into the tank wall at the observation level. A DSLR fitted with a $5\times$ microscope objective was aligned with the slit to record flocs as they crossed the narrow optical path, enhancing image clarity and increasing the number of focused particles captured by the camera, following methods similar to \citeA{osborn2021flocarazi}.

\begin{figure}[htbp]
    \centering
    \includegraphics[width=1\textwidth]{figures/fig2.pdf}
    \caption{
    (a) Schematic of the custom flocculation apparatus. A cylindrical PVC tank (20~cm diameter, 60~cm height) is filled to 40~cm depth and stirred with a motor-driven rotating lid just below the surface. An observation slit (3~mm wide) enables in situ imaging of suspended flocs with a DSLR camera and $5\times$ microscope objective aligned for high-resolution focus. (b) Grain-size distributions of the dispersed input sediment measured by the Mastersizer from \citeA{nghiem2025flocculation}. Curves show pure silica sand (SiO$_2$), montmorillonite clay, and their 50:50 volume-weighted mixture.}
    \label{fig:2}
\end{figure}

\subsection{Experimental Runs}

We conducted eight experimental runs in total. The first two were calibration runs using laboratory spheres and silica grains, which do not flocculate, to verify the accuracy of our settling-velocity measurements and provide a baseline for comparing non-flocculating particles to flocs under identical experimental conditions. The remaining six runs systematically varied shear velocity and organic matter (OM) content, given that previous studies have shown that turbulent shear limits floc size by promoting breakup, while organic matter enhances aggregation by binding particles \cite{nghiem2022mechanistic}. By independently controlling these variables, we aimed to isolate their respective roles in floc formation and settling.

Each floc experiment used 10~g of solids: 5~g silica (SiO$_2$) and 5~g montmorillonite clay. Organic content was varied in three runs by adding guar gum, a polysaccharide proxy for plant-derived extracellular polymers, at 1, 2, or 3~wt\% relative to the total mass of solids, following the organic-loading framework of \citeA{zeichner2021early}. This range was selected to (i) remain within the lower end of organic fractions reported for natural river suspended sediment (order 1–6\% and up to $\sim$20\% in compiled field datasets), and (ii) avoid polymer-dominated regimes while still producing a resolvable change in flocculation dynamics. In \citeA{zeichner2021early}, polymer-to-clay ratios were chosen explicitly to span modern river organic-matter concentrations, and their experiments show that even 1–2~wt\% guar gum produces a strong flocculation response. Guar gum is a controlled proxy for natural extracellular polymeric substances and does not capture the full chemical complexity of field organic matter. Full experimental conditions, including organic-matter fraction and imposed $u_*$, are listed in Table~\ref{tab:exp_conditions}.

\begin{table}[htbp]
\caption{Summary of experimental conditions.\label{tab:exp_conditions}}
\centering
\begin{tabular}{lccc}
\hline
Exp Num & Experiment Type & OM (wt\%) & Shear velocity (m/s) \\
\hline
1 & Lab Sphere     & NA    & NA    \\
2 & Silica         & NA    & NA    \\
3 & Floc Shear 1   & 2     & 0.02  \\
4 & Floc Shear 2   & 2     & 0.03  \\
5 & Floc Shear 3   & 2     & 0.04  \\
6 & Floc OM 1     & 1     & 0.04  \\
7 & Floc OM 2     & 2     & 0.04  \\
8 & Floc OM 3     & 3     & 0.04  \\
\hline
\multicolumn{4}{l}{NA: Not applicable.}
\end{tabular}
\end{table}

Grain-size distributions for the silica and montmorillonite were measured using a Malvern Mastersizer 3000 (Fig.~\ref{fig:2}b), following the procedure of \citeA{nghiem2025flocculation}. In brief, samples were dispersed in deionized water, sonicated to break up aggregates, and analyzed by laser diffraction to obtain particle-size distributions. Together, these two minerals span nearly three orders of magnitude in diameter, with primary particles ranging from $2 \times 10^{-7}$ to $2 \times 10^{-4}$~m (percent finer by volume). For the mixed-sediment runs, a composite distribution was calculated as the 50:50 volume-weighted average of the two minerals.

\subsection{Image Processing}

The floc tracking analysis began with processing every frame of the recorded video sequence, acquired at a constant frame interval $\Delta t = 1/30$~s, as read from the video metadata, so that the sampling frequency equals the acquisition rate. Each frame was converted to an eight-bit grayscale image using OpenCV, and a time-averaged background image (computed by averaging a block of 1000 consecutive frames) was subtracted from each image to isolate transient features such as moving flocs (Fig.~\ref{fig:3}a).

\begin{figure}[htbp]
    \centering
    \includegraphics[width=1\textwidth]{figures/fig3.pdf}
    \caption{
    Image processing workflow used for floc analysis. (a) Example of a background-subtracted image, highlighting transient objects such as flocs. (b) Detected and contoured particles passing focus and sharpness thresholds (outlined in red) with raw velocity vectors overlaid. (c) Local background flow field calculated by PIV from tracer particles. (d) Example of a tracked floc with both its raw and corrected (background-subtracted) settling velocity vectors indicated.
    }
    \label{fig:3}
\end{figure}

Particle boundaries in each image were detected from intensity gradients using a global threshold of 200 applied to the normalized residual image. To retain only sharp and well-focused flocs, candidate objects were required to exceed a Laplacian variance threshold of 5 (measuring focus/sharpness) and to have an intensity standard deviation greater than 5 within their contour \cite{pertuz2013analysis}. These specific threshold values were optimized for our experimental conditions and image resolution and may vary between different setups. Only particles meeting both criteria were accepted for further analysis and are highlighted in red in Fig.~\ref{fig:3}b. This objective, per-frame filtering mitigates focus-related sampling bias by excluding blurred detections rather than allowing size-dependent blur to enter the floc statistics. For each accepted object, we recorded the planform area in image pixels $A_{\mathrm{px}}$ (later converted to an equivalent diameter after spatial calibration) and the centroid coordinates $(x_c,y_c)$ (for subsequent velocity estimation).

To express diameters and velocities in physical units, we calibrated pixel size using a precision stage micrometer imaged under the same optical conditions as the experiments. Comparing the pixel spacing between micrometer markings with their true separations yielded a scale factor of \(0.45~\mu\mathrm{m\,px^{-1}}\). We applied this factor to all subsequent analyses to convert particle dimensions and inter-frame displacements from pixels to micrometers; velocities were then obtained from these displacements using the recorded frame interval.

Settling velocities were obtained by frame-to-frame tracking using a conservative nearest-neighbor association. For each detection in frame \(t\), we searched within a 50-pixel radius in frame \(t{+}1\) for candidate matches. We scored each candidate \(j\) using two equally weighted, normalized differences: the equivalent concave-hull diameter \(d_{\mathrm{ch}}\) (px) and the Laplacian-variance sharpness \(L\) (arb. units). The score was defined as \(S_j=\tfrac{1}{2}\big(|d_{\mathrm{ch},j}-d_{\mathrm{ch},t}|/d_{\mathrm{ch},t}+|L_j-L_t|/L_t\big)\). We selected the candidate with the smallest \(S_j\), provided that both relative differences were \(\le 0.10\); otherwise, the track was terminated to avoid spurious links from rapid shape or focus changes. The vertical velocity for that interval was initially computed as \(\Delta y\,s/\Delta t\), where \(\Delta y\) is the vertical centroid shift (px, positive downward), \(s\) is the spatial calibration (\(0.45~\mu\mathrm{m\,px^{-1}}\)), and \(\Delta t\) is the frame interval (s). Raw velocities computed in \(\mu\mathrm{m\,s^{-1}}\) were subsequently converted to \(\mathrm{m\,s^{-1}}\) to match the physical units required by Eq.~\eqref{eq:stokes} and Eq.~\eqref{eq:floc_settling}.

Frame-to-frame particle displacements yield raw 2D velocity vectors that combine gravitational settling with background flow in the imaging plane. To isolate the settling signal, we first computed a local 2D flow field using OpenPIV \cite{liberzon2021openpiv} for each frame pair. Following \citeA{smith2015image}, the images were digitally filtered to retain only tracers smaller than $10~\mu\mathrm{m}$ that behave as passive flow followers (Fig.~\ref{fig:3}c). For each tracked floc, we then subtracted the corresponding local flow vector---bilinearly interpolated at the floc centroid---from the floc's raw trajectory vector. This vector differencing produces a background-corrected motion (Fig.~\ref{fig:3}d, gray arrow for raw, white arrow for corrected), from which we report the vertical component as the settling velocity \(w_s\). All settling velocities in this study refer to this corrected, background-subtracted quantity.

After velocity correction, floc diameters were derived from each particle's projected area. Two common image-based metrics are used to estimate diameter: the Feret diameter (maximum caliper length across the projection) and the equivalent circular diameter (ECD), defined as the diameter of a circle with the same projected area \cite{jarvis2005measuring, smith2015image}. The Feret diameter is sensitive to orientation and outliers, often overestimating size for irregular or elongated flocs. By contrast, the area-equivalent diameter---especially when computed from a concave hull---can better represent irregular shapes by de-emphasizing extreme protrusions, though it may underrepresent thin filaments. Accordingly, we adopt the concave-hull ECD (with $\alpha = 0.05$, the concave-hull radius parameter that controls the allowed boundary indentation), corresponding to the projected-area--equivalent diameter of \citeA{smith2015image}, while acknowledging that alternative definitions are reasonable \cite{jarvis2005measuring}. Particles were segmented with binary masks, boundaries were traced using a concave-hull algorithm, the enclosed area was measured, and the diameter of a circle of equal area was computed and used as \(d_f\) in all subsequent log--log analyses. Because $d_f$ is defined geometrically from the 2D projection, this step does not assume constant floc density; any density changes associated with size or organic content enter only through the measured settling velocities.

\subsection{Calibration and Settling Analysis}

We calibrated our setup with Experiments~1 and~2, which measured the settling of laboratory spheres (50~\textmu m) and natural silica grains (Fig.~\ref{fig:4}a). Settling velocities were fitted with Stokes' law (Eq.~\eqref{eq:stokes}), treating the shape factor $b_1$ as a free parameter. For the spheres, the best-fit value $b_1 = 17.7$ is consistent with the theoretical value of $18$ for smooth spheres (Fig.~\ref{fig:4}b), which lies within the confidence interval of the fit. This agreement indicates that the experimental setup accurately reproduces Stokes' settling behavior for smooth spheres. In contrast, silica grains settled more slowly, yielding $b_1 = 29.2$, a drag increase attributable to their angular, non-spherical shapes; \citeA{ferguson2004simple} reported a comparable $b_1 \approx 24$ for natural sediments of similar morphology.

Compared with silica, flocs settled more slowly and showed substantially greater dispersion in velocity at a given diameter (Fig.~\ref{fig:4}b). This reduced settling reflects the lower effective density of flocs arising from their highly porous, irregular internal structure \cite{kranenburg1994fractal, winterwerp1998simple}. Consequently, the floc data are not described by the standard Stokes formulation (Eq.~\eqref{eq:stokes}) and instead require the floc-specific relation (Eq.~\eqref{eq:floc_settling}), which introduces the 3D fractal dimension $n_f$ and the primary particle diameter $d_p$ as additional parameters.

\begin{figure}[htbp]
    \centering
    \includegraphics[width=1\textwidth]{figures/fig4.pdf}
    \caption{
    (a) Images of a laboratory sphere (Exp.~1), silica grain (Exp.~2), and a representative floc (Exp.~5). The 30~$\mu$m scale bar applies to all. (b) Log--log plot of settling velocity ($w_s$) vs.\ diameter ($d_f$) for spheres (green circles, Exp.~1), silica (red circles, Exp.~2), and aggregated flocs (blue circles, Exps.~3--8). Sphere and silica fits use Stokes' law (Eq.~\eqref{eq:stokes}), yielding shape factors ($b_1$) of 17.7 and 29.2. The floc data fit the fractal settling model ($w_s \propto d_f^{n_f - 1}$) with a slope of 0.90, giving an apparent fractal dimension $n_f = 1.90$. (c) Sensitivity of the fractal settling model (Eq.~\eqref{eq:floc_settling}) to assumed primary particle diameter ($d_p$). The floc data are compared to theoretical predictions using $n_f = 1.9$ for a primary particle size ($d_p$) of $10^{-6}$~m (dashed line) and $10^{-5}$~m (solid red line). This highlights how the choice of $d_p$ controls the required model fit, a key uncertainty in the conventional approach.
    }
    \label{fig:4}
\end{figure}

However, directly calculating $n_f$ by fitting Eq.~\eqref{eq:floc_settling} to the floc data is problematic because the primary-particle diameter, $d_p$, is unknown. Although our experiments have known distributions of input particle sizes (Fig.~\ref{fig:2}b), we do not know which particle sizes are actually incorporated into any given floc. Even if we could measure the primary particle sizes directly, the fractal theory necessitates specifying only a single value for a characteristic primary particle size, and it is unclear how to define this value from a distribution. Moreover, because $n_f$ appears both as an exponent on $d_f$ and within the pre-factor of Eq.~\eqref{eq:floc_settling} (via $d_p^{-(n_f-3)}$), a conventional regression that fits the model to the data for a prescribed $d_p$ (a one-parameter fit) forces the inferred $n_f$ to mathematically compensate for the assumed pre-factor. Consequently, the assumed value for $d_p$ exerts a strong influence on the inferred fractal dimension $n_f$ (Fig.~\ref{fig:4}c): assuming $d_p = 1\times 10^{-6}$~m yields $n_f = 2.77$; $d_p = 1\times 10^{-5}$~m gives $n_f = 2.10$; and $d_p = 1\times 10^{-4}$~m produces $n_f = 3.00$.

To avoid explicitly solving for the unknown primary-particle size $d_p$, we use the logarithmic form of Eq.~\eqref{eq:floc_settling}: $\log(w_s) = (n_f - 1)\log(d_f) + C$. Here, the intercept $C$ lumps together $d_p$, the shape factor $b_1$, and other constants. Under the assumption that \(d_p\) and \(b_1\) are approximately constant across the fitted population, the slope in \(\log(w_s)\)--\(\log(d_f)\) space yields \(n_f = \text{slope} + 1\) without requiring prior knowledge of \(d_p\). For the full floc population, the fitted slope is \(0.90\), giving \(n_f = 1.90\) (Fig.~\ref{fig:4}b). However, if structural heterogeneity causes \(n_f\) to covary with size, or if effective \(d_p\) varies across the population, a bulk regression can exhibit Simpson-type data-aggregation bias. We therefore stratify flocs by their in situ, image-based fractal dimension and fit separate regressions within each stratum; the resulting slopes provide subgroup-specific estimates of \(n_f\) and reveal systematic dependence on structure.

\subsection{Image-based 2D box-counting dimension}

Fractal theory relates the fractal dimension to the primary particle diameter, the number of primary particles, and the floc diameter (Eq.~\eqref{eq:fractaltheory}). In natural systems, however, neither $d_p$ nor $N$ is constant or directly measurable during experiments. As a result, numerous methods have been developed to estimate floc fractal dimensions from image-based properties. These approaches are inherently 2D, yielding an apparent fractal dimension that may differ from the hydrodynamically relevant 3D value. Throughout this paper, we denote the 3D fractal dimension inferred from settling velocity scaling as $n_f$, and use $n_f^{(\mathrm{2D})}$ for the 2D box-counting dimension measured from images. Where the specific 2D method matters (Methodology), we distinguish the box-counting dimension $n_{f,\,\mathrm{box}}^{(\mathrm{2D})}$ from the perimeter--area dimension $n_{f,\,\mathrm{perimeter}}^{(\mathrm{2D})}$; in the Results, all 2D values refer to box-counting and are written simply as $n_f^{(\mathrm{2D})}$.

The perimeter--area method has previously been used to estimate a 2D fractal dimension of flocs~\cite{lee2004prediction, maggi2004method, tang2015reconstructing}. This method relies on the empirical scaling relation:
\begin{equation}
P \propto A^{n_{f,\,\mathrm{perimeter}}^{(\mathrm{2D})} / 2}
\end{equation}
where $P$ is the measured perimeter and $A$ is the projected area of the particle in the 2D image. In practice, $n_{f,\,\mathrm{perimeter}}^{(\mathrm{2D})}$ is estimated from the log--log slope of $P$ versus $A$. This expression is derived by analogy to the Euclidean relation $P = k A^{1/2}$, which holds for smooth (non-fractal) boundaries with $n_{f,\,\mathrm{perimeter}}^{(\mathrm{2D})} = 1$. The underlying assumption is that, for a fractal boundary, the exponent in this scaling law approximates the true boundary fractal dimension $n_{f,\,\mathrm{perimeter}}^{(\mathrm{2D})}$. However, both mathematical analysis and geometric constructions (e.g., the Koch curve) demonstrate that this assumption breaks down for truly fractal sets~\cite{cheng1995perimeter,frame_mandelbrot_neger_2025}. For such boundaries, the measured perimeter is strongly dependent on the image resolution, making a simple $P$--$A$ scaling ill-defined without explicitly accounting for the measurement scale. As a result, perimeter--area scaling does not reliably recover the true fractal dimension from 2D projections.

Box-counting, in contrast, is an appropriate method for fractal-dimension estimation that does not rely on strict self-similarity and is computationally straightforward \cite{foroutan1999advances}. A key advantage of box-counting is that it does not require knowledge of the primary particle size ($d_p$) or the number of constituent particles ($N$), making it particularly well-suited for natural flocs where these values are unknown or variable~\cite{bellouti1997flocs, spencer2022quantification}. The box-counting dimension $n_{f,\,\mathrm{box}}^{(\mathrm{2D})}$ is determined by measuring how the number of occupied boxes $N_{\text{box}}(\epsilon)$ at a given scale $\epsilon$ varies with scale~\cite{mandelbrot1967long}:
\begin{equation}
N_{\text{box}}(\epsilon) \propto \epsilon^{-n_{f,\,\mathrm{box}}^{(\mathrm{2D})}}
\end{equation}

For each floc mask, we computed $N_{\text{box}}(\epsilon)$ using square boxes with side lengths $\epsilon$ spanning powers of two in pixel units (\(\epsilon=1,2,4,\dots\)), and we estimated the dimension from the slope of a least-squares fit over the portion of the log--log curve that remains approximately linear and spans at least one order of magnitude in scale. When analyzing images of objects with finite extent from in situ imaging, finite-size effects and image pixelation make the estimated fractal dimension sensitive to the available scaling range. For a digital image of $W \times W$ pixels, the practical range of box sizes ($\epsilon$) spans from 1 to $W$ pixels. Reliable estimation of $n_{f,\,\mathrm{box}}^{(\mathrm{2D})}$ requires a sufficiently broad range of box sizes to establish a stable linear trend on a log--log plot of $N_{\text{box}}(\epsilon)$ versus $1/\epsilon$. If the image is too small, the limited scaling range can yield unstable values of $n_{f,\,\mathrm{box}}^{(\mathrm{2D})}$. For example, a $20 \times 20$ pixel image provides only about 1.3 orders of magnitude in scale, which is insufficient for robust fitting and produces scale-dependent estimates (Fig.~\ref{fig:5}a). In contrast, a $100 \times 100$ pixel image spans two full orders of magnitude and typically provides 6--7 usable data points for regression; beyond this image size, the estimated $n_{f,\,\mathrm{box}}^{(\mathrm{2D})}$ saturates, indicating a practical threshold for resolution-independent estimation (Fig.~\ref{fig:5}a). Accordingly, all analyzed flocs were extracted using $100 \times 100$ pixel windows centered on each particle. Fig.~\ref{fig:5}b shows examples for two experimental flocs at high image resolution, illustrating the application of the box-counting method to real data. In this study, all reported values of the 2D box-counting dimension ($n_{f,\,\mathrm{box}}^{(\mathrm{2D})}$) in the results and discussion sections refer exclusively to measurements obtained by the box-counting method.

\begin{figure}[htbp]
    \centering
    \includegraphics[width=1\textwidth]{figures/fig5.pdf}
    \caption{
    (a) Illustration of resolution dependence in box-counting dimension estimation. Top: A perfect circle embedded in $20 \times 20$ and $200 \times 200$ pixel grids shows that coarse grids systematically underestimate the fractal dimension. Bottom: Plot of measured 2D box-counting dimension $n_{f,\,\mathrm{box}}^{(\mathrm{2D})}$ versus image width $W$, highlighting a resolution-dependent regime at low $W$ and a stable region at high $W$. The transition scale ($W \approx 100$ pixels) was identified by repeating the box-counting calculation across increasing image sizes and selecting the smallest $W$ beyond which further increases produced negligible changes in the estimated dimension.
    (b) Example box-counting dimension estimates for two experimental flocs imaged at $>100 \times 100$ pixel resolution.
    }
    \label{fig:5}
\end{figure}

\subsection{Primary Particle Diameter ($d_p$)}

In most in situ experiments, the primary particle diameter \(d_p\) is assumed from the input sediment-size distribution and hypothesized floc-formation conditions because \(d_p\) is rarely resolvable directly. This practice can be uncertain when the effective primary scale differs from the feed distribution or varies across aggregation pathways. Here we identify \(d_p\) empirically by leveraging multiple floc subpopulations with distinct fractal dimensions \(n_f\). In the experimental application, we treat $d_p$ as an experiment-level parameter common to all $n_f$-stratified subpopulations. As implied by Eq.~\eqref{eq:floc_settling}, \(d_p\) does not affect the slope of the \(\log(w_s)\)--\(\log(d_f)\) relation (the slope is \(n_f-1\)), but it fixes a common intersection of the subgroup regressions. This is because at \(d_f=d_p\), the object is a primary particle (not a floc), so substituting \(d_f=d_p\) into Eq.~\eqref{eq:floc_settling} yields \(w_s=(g\,R_s/b_1\nu)\,d_p^{2}\), which is independent of \(n_f\). Consequently, regressions from any number of \(n_f\)-stratified subpopulations intersect at \((d_f, w_s)=(d_p,\,(g\,R_s/b_1\nu)\,d_p^{2})\). We refer to this common intersection diameter as $d_\text{intersect}$. In our synthetic example (Fig.~\ref{fig:1}), three subgroups with \(n_f \approx 1.5,\,2.0,\) and \(2.5\) produce regression lines that converge at \(d_\text{intersect} \approx 10^{-5}\,\mathrm{m}\), because the synthetic data were generated using this prescribed primary-particle diameter across varying fractal dimensions, thereby illustrating the method.


A floc contains many primary particles. For example, with $d_f \sim 10^{-4}$~m, $d_p \sim 10^{-5}$~m, and $n_f \approx 2$, Eq.~\eqref{eq:fractaltheory} gives $N \sim (d_f/d_p)^{n_f} \sim 10^{2}$–$10^{4}$. This means the settling measurements reflect the collective behavior of aggregates built from a common sediment pool, and the inferred $d_p$ should be interpreted as an effective primary scale for the flocs resolved by the imaging and tracking. In the synthetic example, the subgroup regressions intersect exactly at a single point because the data are generated under the idealized assumption that all primary particles share a single diameter $d_p$. In natural sediments, primary particles span a distribution of sizes, so different flocs can be built from differently sized particles and the subgroup regressions need not converge to a unique intersection. We therefore treat the intersection as distributed rather than singular and use a Monte Carlo procedure to generate an ensemble of intersection points, which we interpret as a distribution of effective $d_p$ values. We then compare this inferred $d_p$ distribution to the input sediment-size distribution to assess consistency and to quantify how the sediment’s primary-size variability propagates into the settling-based estimate of $d_p$.

The intersection method assumes a common shape factor $b_1$ across fractal subgroups. Because drag depends on aggregate structure, $b_1$ may vary with $n_f$, biasing the inferred intersection diameter. Allowing for subgroup-dependent drag, the relationship between the observed intersection diameter $d_\text{intersect}$ and the true primary particle diameter $d_p$ for two subpopulations $j$ and $k$ is
\begin{equation}
d_\text{intersect} = d_p \left( \frac{b_{1,j}}{b_{1,k}} \right)^{1/(n_{f,j} - n_{f,k})},
\label{eq:dpintersection}
\end{equation}
which reduces to
\begin{equation}
d_\text{intersect} = d_p
\label{eq:dpintersectionideal}
\end{equation}
when $b_{1,j}=b_{1,k}$. In this limit, all subgroup regressions intersect exactly at the primary particle diameter.

\section{Results}

\subsection{Floc Bulk Properties}
\label{sec:bulk}
We first report the bulk distributions of floc diameter ($d_f$) and settling velocity ($w_s$) across all experimental conditions (Fig.~\ref{fig:6}). Median $d_f$ decreases monotonically with increasing shear velocity (Fig.~\ref{fig:6}a), consistent with stronger turbulence breaking apart aggregates. In contrast, $d_f$ responds weakly and non-monotonically to organic matter content, increasing slightly from 1\% to 2\% OM before declining at 3\% OM (Fig.~\ref{fig:6}b). Median settling velocity $w_s$ exhibits a non-monotonic dependence on shear velocity, peaking at $u_* = 0.03$~m\,s$^{-1}$ rather than decreasing steadily (Fig.~\ref{fig:6}c). With increasing organic matter, $w_s$ decreases monotonically (Fig.~\ref{fig:6}d).

\begin{figure}[htbp]
    \centering
    \includegraphics[width=0.8\textwidth]{figures/fig6.pdf}
    \caption{Median floc diameter ($d_f$) and settling velocity ($w_s$) with interquartile ranges for varying shear velocity (a, c) and organic matter content (b, d). Dashed lines show normalized model predictions: in panels (a, b), the \citeA{nghiem2022mechanistic} equilibrium floc diameter trend $d_{f,\text{eq}} = f(\eta,\theta)$, which is independent of $n_f$; in panels (c, d), the base model with $n_f = 2$ (red dashed; Eq.~\eqref{eq:floc_settling} with $n_f=2$) and the modified model with measured $n_f$ and $d_p$ (blue dashed; Eq.~\eqref{eq:floc_settling} with condition-specific $n_f$ and $d_p$; see Section~\ref{sec:reconciling}). Each model trend is normalized by a single least-squares constant.}
    \label{fig:6}
\end{figure}


To interpret these trends mechanistically, we compare our observations against the semi-empirical freshwater flocculation framework of \citeA{nghiem2022mechanistic}. That model predicts both the equilibrium floc diameter $d_f$ and the corresponding settling velocity $w_s$ as functions of two environmental inputs: the Kolmogorov microscale $\eta$ (which sets the smallest turbulent eddy that can break flocs) and the fractional organic surface coverage $\theta$ (which controls binding strength). Evaluating this model against our data requires mapping our experimental control variables---imposed shear velocity $u_*$ and organic matter weight fraction---to $\eta$ and $\theta$.

In applying their model to river field data, \citeA{nghiem2022mechanistic} used an open-channel relation where turbulent dissipation is generated solely by bed friction. Our experimental apparatus is a top-driven cylindrical tank in which stationary sidewall friction acts as an additional dominant energy sink. To calculate the volume-averaged turbulent dissipation rate $\bar{\varepsilon}_{\text{tank}}$, we evaluate the steady-state power balance \cite{camp1943velocity}, equating the power injected by the rotating lid to the power dissipated by friction on all tank boundaries (bed and sidewalls). Assuming the wall shear stress scales with the calibrated bed shear velocity ($\tau_w \approx \rho_w u_*^2$) \cite{Pope2000, Schlichting2000} and substituting the bulk slip velocity $U_{\text{bulk}} \approx u_* \sqrt{2/C_f}$, the mean dissipation rate is:
\begin{equation}
    \bar{\varepsilon}_{\text{tank}} = \frac{u_*^3}{\sqrt{C_f/2}} \left( \frac{1}{h} + \frac{2}{R} \right),
    \label{eq:eps_tank}
\end{equation}
where $h = 0.40$~m is the water depth, $R = 0.10$~m is the tank radius, and $C_f \approx 0.005$ is the skin friction coefficient for fully turbulent swirling flow \cite{daily1960chamber, Schlichting2000}. The Kolmogorov microscale is then $\eta = (\nu^3 / \bar{\varepsilon}_{\text{tank}})^{1/4}$. Because of the extensive sidewall area, $\bar{\varepsilon}_{\text{tank}}$ is substantially larger than open-channel predictions for an equivalent $u_*$, so the absolute value of $\eta$ is smaller. However, $\bar{\varepsilon}_{\text{tank}} \propto u_*^3$ regardless of tank geometry, so the functional scaling is identical to an open channel: $\eta \propto u_*^{-3/4}$.

\citeA{nghiem2022mechanistic} parameterized organic binding using field measurements of particulate organic carbon (OC), requiring an intermediate stoichiometric assumption (cellulose composition) to convert OC to total organic matter (OM). Because our experiments use guar gum---a pure polysaccharide proxy for natural extracellular polymeric substances---we bypass this conversion and use the experimental weight percentage of OM directly. Following Eq.~19 of \citeA{nghiem2022mechanistic}, the fractional surface coverage is:
\begin{equation}
    \theta = \frac{(\text{\%OM}/100)\,(\rho_s / \rho_{\text{OM}})\, d_p^3}{(d_p + \delta)^3 - d_p^3},
    \label{eq:theta}
\end{equation}
where $\rho_s = 2650$~kg\,m$^{-3}$ is the sediment grain density, $\rho_{\text{OM}} = 1000$~kg\,m$^{-3}$ is the organic matter density, $d_p$ is the primary particle diameter, and $\delta = 1~\mu$m is the assumed organic coating thickness.

A key property of the \citeA{nghiem2022mechanistic} framework is that the predicted equilibrium floc diameter does not depend on the fractal dimension. In their model (Eq.~9 in their study), the fractal dimension $n_f$ enters both the aggregation and breakup terms of the floc growth equation through identical powers of $d_f$: aggregation scales as $d_f^{\,n_f+2}$ and breakup as $d_f^{\,n_f}$. When these two terms are set equal at equilibrium, $n_f$ cancels algebraically and the resulting equilibrium diameter $d_{f,\text{eq}}$ depends only on $\eta$ and $\theta$. Physically, this cancellation reflects the fact that the same fractal structure governs both how efficiently particles collide (aggregation) and how susceptible the aggregate is to turbulent stresses (breakup), so the two effects offset exactly at equilibrium.

Using the calibrated exponents from \citeA{nghiem2022mechanistic} (their Eqs.~23--24), the equilibrium diameter scales as $d_{f,\text{eq}} \propto \eta^{r}[\theta^2(1-\theta)^2]^{q}$, where $r$ and $q$ are fitted constants. Since $\eta \propto u_*^{-3/4}$, the predicted shear dependence is $d_{f,\text{eq}} \propto u_*^{-0.75r}$, and the organic-matter dependence enters through $\theta$ (Eq.~\eqref{eq:theta}). Because the model parameters were calibrated for river field conditions rather than our tank geometry, we compare trends rather than absolute magnitudes: we normalize the predicted $d_{f,\text{eq}}$ by a single least-squares constant and plot it against the measured medians in Fig.~\ref{fig:6}a,b; the model captures the monotonic decrease of $d_f$ with shear and the weak non-monotonic response to organic matter.

Although the equilibrium floc diameter is independent of $n_f$, the settling velocity is not. Equation~\eqref{eq:floc_settling} shows that $w_s$ depends on $n_f$ through two exponents: $d_p^{\,3-n_f}$ in the prefactor and $d_f^{\,n_f-1}$ in the size scaling. The \citeA{nghiem2022mechanistic} model adopts $n_f = 2$ throughout, which simplifies Eq.~\eqref{eq:floc_settling} to $w_s \propto d_p\, d_f$, a linear dependence on floc diameter in which $d_p$ and $b_1$ are absorbed into the prefactor. We refer to this as the base model. We compare the predicted and observed trends rather than absolute magnitudes for two reasons. First, the \citeA{nghiem2022mechanistic} model parameters were calibrated against river field data; applying them to our tank geometry, which has a different energy dissipation distribution, would require re-calibration. Second, the settling velocity prefactor contains the shape factor $b_1$, which is not independently constrained and may itself vary with $n_f$ (see Section~\ref{sec:dp}). We therefore normalize each model prediction by a single least-squares constant ($\hat{w}_{s,\text{base}} = k_{\text{base}}\, d_{f,\text{eq}}$) and compare the predicted trend against the measured medians in Fig.~\ref{fig:6}c,d. Because the base model is linear in $d_f$, and $d_f$ decreases monotonically with shear, the base model strictly predicts a monotonic decrease in $w_s$ with increasing $u_*$. It therefore cannot reproduce the non-monotonic peak observed at $u_* = 0.03$~m\,s$^{-1}$ (Fig.~\ref{fig:6}c). Similarly, because $d_f$ increases slightly with OM (for $\theta < 0.5$), the base model predicts $w_s$ increasing with OM, opposite to the observed monotonic decline (Fig.~\ref{fig:6}d).

These discrepancies motivate a closer examination of the internal structural properties---the fractal dimension $n_f$ and the effective primary particle diameter $d_p$---that the base model holds fixed. In the following two subsections, we develop methods to extract condition-specific estimates of $n_f$ (Section~\ref{sec:stratifying}) and $d_p$ (Section~\ref{sec:dp}), and in Section~\ref{sec:reconciling} we show that incorporating this measured structural variability into the settling velocity prediction resolves the failures identified above.


\subsection{Stratifying 2D Box-Counting Data}
\label{sec:stratifying}

To infer a fractal dimension from settling data, we group flocs into bins based on their 2D box-counting dimension, $n_f^{(\mathrm{2D})}$, measured from images. The bin width in $n_f^{(\mathrm{2D})}$ serves only as a coarse-graining scale for separating discrete structural classes, not as a physical control parameter. If a population contains distinct structural classes (e.g., non-flocculating grains with $n_f \approx 3$ and flocs with $n_f \approx 2$), finer stratification merely subdivides these classes into redundant bins that recover the same class-specific slopes. 

As a sensitivity analysis, we computed the 2D box-counting dimension, $n_f^{(\mathrm{2D})}$, for Experiment~3 (2~wt\% organic matter; shear velocity $u_* = 0.02~\mathrm{m\,s^{-1}}$). To evaluate the impact of stratification, we divided the dataset into bins of width 0.1, 0.05, and 0.025, corresponding to 3, 6, and 12 bins of $n_f^{(\mathrm{2D})}$, respectively, within a manually defined range from 1.3 to 1.6. Within each bin, we fitted a power-law relation between settling velocity, $w_s$, and floc diameter, $d_f$, in log--log space (Figs.~\ref{fig:7}a, d, g). The model $\log_{10}(w_s) = m \log_{10}(d_f) + C$ was calibrated using an ensemble Markov chain Monte Carlo (MCMC) sampler to obtain robust parameter estimates and verify convergence. After burn-in and thinning, the posterior medians of the slope $m$ were converted to the hydrodynamically relevant 3D fractal dimension, denoted simply as $n_f$, using the relation $n_f = m + 1$.

Plotting the inferred 3D fractal dimension, $n_f$, against the corresponding group-mean 2D value, $n_f^{(\mathrm{2D})}$, shows how the stratification behaves across binning schemes. A linear trend captures the binned data reasonably well for all cases (Figs.~\ref{fig:7}c, f, i). However, because the number of bins is small (3, 6, or 12), the resulting coefficients of determination are not themselves a robust metric of model adequacy, as high $R^2$ values are expected even in the presence of substantial uncertainty.

A more informative assessment comes from quantifying the dynamical consequences of refining the stratification. From Eq.~\eqref{eq:floc_settling}, assuming the log--log regression lines converge near the primary particle diameter $d_p$, a change $\Delta n_f$ produces a vertical shift $\Delta \log_{10}(w_s) \approx \Delta n_f \log_{10}(d_f/d_p)$. Over the observed range of macroscopic floc sizes ($d_f/d_p \approx 10^1$--$10^2$), a change of $\Delta n_f = 0.01$ produces a shift of $\Delta \log_{10}(w_s) \approx 0.01$--$0.02$, corresponding to a change in settling velocity of less than 5\%. This effect is minimal relative to measurement scatter. Consequently, increasingly fine stratification does not reveal qualitatively different settling regimes but instead primarily inflates variance in the estimated slopes as bin populations decrease. We therefore adopt a bin width of 0.1 for subsequent analyses. This choice resolves meaningful structural heterogeneity while maintaining sufficient sample sizes for regression, and avoids over-resolving differences that have negligible impact on settling velocities over the observed floc-size range.

\begin{figure}[htbp]
    \centering
    \includegraphics[width=1\textwidth]{figures/fig7.pdf}
    \caption{Stratification of floc data by 2D box-counting dimension (\(n_f^{(\mathrm{2D})}\)) for Experiment~3, illustrating how binning strategy influences the estimation of 3D fractal dimensions. (a--c) Analysis using 3 bins: (a) log--log plots of settling velocity (\(w_s\)) versus floc diameter (\(d_f\)), color-coded by 2D box-counting dimension bin; (b) fitted slopes (\(m\)) for each bin; (c) resulting 3D fractal dimension (\(n_f = m + 1\)) plotted against mean \(n_f^{(\mathrm{2D})}\) for each bin, showing a strong linear relationship. (d--f) Same analysis using 6 bins; (g--i) using 12 bins, highlighting increased scatter and reduced regression stability with finer binning.}
    \label{fig:7}
\end{figure}

The six floc experiments (Experiments 3--8) comprise approximately 12{,}500 individual floc particles. Within each 0.1-wide $n_f^{(\mathrm{2D})}$ bin, outliers were removed using the interquartile range criterion [$Q_1 - 1.5\,\mathrm{IQR},\, Q_3 + 1.5\,\mathrm{IQR}$], yielding five to eight bins per experiment with a median population of 190 flocs per bin. The same MCMC regression procedure described above was then applied within each bin to obtain $n_f$.

Figure~\ref{fig:8} shows the log--log relationship between settling velocity $w_s$ and floc diameter $d_f$ for experiments with increasing shear velocity (Figs.~\ref{fig:8}a--c) and increasing organic matter content (Figs.~\ref{fig:8}d--f). Within each panel, power-law regressions fitted to individual $n_f^{(\mathrm{2D})}$ bins reveal distinct slopes that vary systematically with floc structure, and the corresponding 3D fractal dimensions $n_f$ are obtained via the relation $n_f = m + 1$. As the $n_f^{(\mathrm{2D})}$ bin midpoint increases, the recovered $n_f$ also increases monotonically, demonstrating that the 2D structural signature can be used to resolve distinct 3D fractal dimensions within a mixed population. The inset panels illustrate this trend, plotting the bin-median $n_f$ against the $n_f^{(\mathrm{2D})}$ bin midpoint with 95\% credible intervals derived from the MCMC posteriors. For comparison, the aggregated (unstratified) regression and its 95\% credible interval are shown as the orange dashed line and shaded band, respectively. Notably, the aggregated credible interval does not capture the range of $n_f$ values captured by the stratified bins, indicating that a single power-law regression applied to the bulk population cannot capture the structural heterogeneity present within each experiment.

\begin{figure}[htbp]
    \centering
    \includegraphics[width=1\textwidth]{figures/fig8.pdf}
    \caption{Comparison of stratified and aggregated (bulk) fits for all floc datasets. (a--c) Varying shear velocities of 0.02, 0.03, and 0.04~m\,s$^{-1}$. (d--f) Varying organic matter contents of 1\%, 2\%, and 3\%. Inset plots show the relationship between \(n_f^{(\mathrm{2D})}\) and slope-inferred \(n_f\) for each experiment; points represent bin medians (with 95\% credible intervals) and are color-coded by \(n_f^{(\mathrm{2D})}\) group to make within-population variability visually explicit.}
    \label{fig:8}
\end{figure}


The dependence of the stratified and aggregated $n_f$ estimates on shear velocity and organic matter is summarized in Fig.~\ref{fig:9}a,b. In the stratified analysis, the interquartile range (IQR) is computed directly from the population of inferred $n_f$ values across bins, weighted by the number of particles contributing to each bin, because each bin represents an independent subset of flocs with distinct fractal structure. In contrast, in the aggregated analysis the IQR is derived from the Bayesian posterior distribution of the fitted slope parameter, since all particles are pooled into a single regression and variability arises from parameter uncertainty rather than particle-to-particle heterogeneity.

Stratified median $n_f$ increases approximately linearly with shear velocity at a rate of $\Delta n_f / \Delta u_* \approx 17$~s\,m$^{-1}$ ($n_f = 17.4\,u_* + 1.42$, Fig.~\ref{fig:9}a), consistent with shear-driven breakup preferentially destroying loose, low-$n_f$ flocs and favoring compact, high-$n_f$ flocs. The aggregated analysis yields a nearly flat response ($n_f = 1.0\,u_* + 1.94$), demonstrating that the shear-driven structural trend is effectively invisible in bulk data. For organic matter (Fig.~\ref{fig:9}b), stratified median $n_f$ decreases at $\Delta n_f / \Delta \mathrm{OM} \approx -0.11$ per wt\% ($n_f = -0.11\,\mathrm{OM} + 1.92$), indicating that organic binding promotes open, porous structures. The aggregated trend is again much weaker ($n_f = -0.03\,\mathrm{OM} + 1.81$). This insensitivity arises because the aggregated regression fits a single power law through a heterogeneous mixture of floc structures, each with a different $w_s$--$d_f$ slope. The resulting best-fit slope is dominated by the most populated $n_f^{(\mathrm{2D})}$ bins and by the large scatter inherent in pooling structurally distinct subpopulations, which masks the systematic variation that the stratified analysis resolves. In effect, the noise introduced by collapsing all fractal classes into one regression absorbs the signal that distinguishes compact from porous flocs.

\begin{figure}[htbp]
    \centering
    \includegraphics[width=1\textwidth]{figures/fig9.pdf}
    \caption{Comparison of stratified and aggregated (bulk) analyses of floc fractal dimensions. (a) Median inferred \(n_f\) versus shear velocity. (b) Median \(n_f\) versus organic matter (OM) content. Error bars denote interquartile ranges (IQRs) and markers indicate medians; for stratified estimates (orange) the IQR reflects particle-scale variability across bins weighted by population, whereas for aggregated estimates (navy) it reflects uncertainty in the Bayesian posterior of a single fitted parameter. (c) Mapping between 2D and 3D fractal dimensions: black dashed line, theoretical relation \(n_f = n_f^{(\mathrm{2D})} + 1\); orange solid line, empirical 2D box-counting to 3D box-counting (BC) conversion ($n_f = 0.81\,(n_f^{(\mathrm{2D})})^{1.81}$); orange dashed line, empirical 2D box-counting to 3D power-law (PL) conversion ($n_f = 0.20\,(n_f^{(\mathrm{2D})})^{4.08}$) \cite{wang2022fractal}. Experimental data (grey markers) align most closely with the empirical box-counting model.}
    \label{fig:9}
\end{figure}

Figure~\ref{fig:9}c collects the per-bin data from all six experiments to examine the mapping between $n_f^{(\mathrm{2D})}$ and $n_f$. The 40 bin-median values (grey circles) are compared against three reference models. The idealized Euclidean relation $n_f = n_f^{(\mathrm{2D})} + 1$ (black dashed) assumes perfect space-filling and systematically overpredicts $n_f$ across the measured range. The empirical power-law (PL) and box-counting (BC) conversions of \citeA{wang2022fractal}, calibrated on computationally generated aggregates, bracket the data more closely. The BC model ($n_f = 0.8118\, {n_f^{(\mathrm{2D})}}^{1.8054}$, solid orange) tracks the upper envelope of the observations, while the PL model ($n_f = 0.2015\, {n_f^{(\mathrm{2D})}}^{4.079}$, dashed orange) underpredicts at moderate $n_f^{(\mathrm{2D})}$. Our data scatter between these two empirical curves, with the closest overall agreement to the BC conversion. This alignment is consistent with the fact that both our stratification variable and the \citeA{wang2022fractal} BC model are based on box-counting dimensions. Mismatch may arise because the \citeA{wang2022fractal} model constructs flocs from monodisperse primary particles of uniform size and shape, whereas natural flocs comprise polydisperse, irregularly shaped constituents whose heterogeneous packing alters the 2D projection geometry.

\subsection{Inferring Primary Particle Diameter ($d_p$)}
\label{sec:dp}

The stratified regressions also provide a means of estimating the effective primary particle diameter, $d_p$. In log--log space, Eq.~\eqref{eq:floc_settling} predicts that regression lines for different $n_f$ converge as $d_f \to d_p$: at the primary-particle scale, a floc reduces to a single grain whose settling velocity is independent of fractal structure. We exploit this convergence by computing all pairwise intersections of the per-bin regression lines from their MCMC posteriors and assume a constant shape factor $b_1$. Repeating the calculation across thousands of posterior draws yields an ensemble of intersection diameters whose spread reflects both measurement scatter and finite-sample uncertainty. Fig.~\ref{fig:10}a--c compare the resulting $d_p$ posterior (orange) with the known input grain-size distribution (gray) for the three shear-velocity experiments, and Fig.~\ref{fig:10}e--g show the corresponding comparison for the three organic-matter experiments. 

\begin{figure}[htbp]
    \centering
    \includegraphics[width=1\textwidth]{figures/fig10.pdf}
    \caption{
        Inferred primary particle diameter ($d_p$) distributions and summary statistics. (a--c) Cumulative percent finer by volume comparing the known input grain-size distribution (grey) with the inferred $d_p$ posterior (orange) for each shear-velocity experiment; KS statistics quantify agreement. (d) Median inferred $d_p$ (orange, with IQR error bars) versus shear velocity, alongside the input-sediment median and IQR (grey). (e--g) Same as (a--c) for the three organic-matter experiments. (h) Median inferred $d_p$ versus organic matter content.
    }
    \label{fig:10}
\end{figure}

Results show agreement between the inferred $d_p$ distribution from flocs and the input grain-size distribution, particularly in their central tendency (e.g., alignment of medians) (Fig.~\ref{fig:10}). The Kolmogorov--Smirnov (KS) statistic ranges from 0.28 to 0.43 across all six experiments, indicating moderate agreement (a KS value of 0 denotes perfect agreement, whereas a value near 1 indicates complete mismatch). Perfect agreement is not expected because the input grain-size distribution in Fig.~\ref{fig:10} includes all particle sizes added to the experiments, but not all of these necessarily serve as primary particles in flocs. In every case the inferred $d_p$ distribution is substantially narrower than the input grain distribution, with posteriors peaking in the coarse-silt range ($\sim$10--40~\textmu m), suggesting that the effective building blocks of the observed flocs are drawn predominantly from the silt fraction \cite{nghiem2026flocculated}.

Panels (d) and (h) of Fig.~\ref{fig:10} summarize the median $d_p$ (orange, with IQR error bars) alongside the input-sediment median and IQR (grey) across shear velocity and organic matter. With increasing shear velocity (panel d), the inferred $d_p$ increases from 21.7~\textmu m at $u_* = 0.02$~m\,s$^{-1}$ to 23.0~\textmu m at 0.03~m\,s$^{-1}$ and 29.8~\textmu m at 0.04~m\,s$^{-1}$, consistent with stronger shear selectively destroying the most fragile aggregates and shifting the effective primary scale toward coarser grains. With increasing organic matter (panel h), the median $d_p$ decreases monotonically from 34.2~\textmu m at 1\% OM to 27.9~\textmu m at 2\% and 17.6~\textmu m at 3\%, indicating that organic binding enables finer particles to participate as effective building blocks. Across all conditions the inferred $d_p$ remains well below the input-sediment median of $\sim$39~\textmu m (grey markers in panels d and h). 

The inferred $d_p$ thus represents an effective primary scale associated with the population of flocs contributing to the measured settling velocities, rather than the full input grain-size distribution. Differences between the input grain-size and inferred $d_p$ distributions reflect both uncertainties in the intersection method and real selective incorporation of certain size classes into flocs. Prolonged turbidity following the experiments indicates that a fraction of very fine material remained in suspension and was not flocculated; however, this observation alone does not distinguish between incomplete aggregation of the smallest particles and limitations of the imaging and tracking workflow at small sizes. Apparent offsets may also arise from size-dependent participation in the floc population resolved by the measurements and from under-sampling of dispersed particles below the detection threshold.

The absolute magnitude of the inferred $d_p$ is sensitive to the assumed value of the shape factor $b_1$. Laboratory and field studies indicate that $b_1$ is not constant and may increase with $n_f$ \cite{strom2011explicit}. If $b_1$ varies from $\mathcal{O}(10)$ to $\mathcal{O}(10^2)$ across fractal classes, the inferred intersection diameter $d_{\mathrm{intersect}}$ can differ from the true $d_p$ by up to an order of magnitude for fixed $n_f$ (Eq.~\eqref{eq:dpintersection}). However, because the same $b_1(n_f)$ relationship applies uniformly across all experiments, the relative trends in $d_p$ with shear velocity and organic matter reported above are preserved regardless of the assumed shape-factor model; only the absolute scale shifts. 

\begin{figure}[htbp]
    \centering
    \includegraphics[width=1\textwidth]{figures/fig11.pdf}
    \caption{
    Influence of the shape factor (\(b_1\)) on the geometric intersection diameter (\(d_\text{intersect}\)) between floc groups of different fractal dimension.
    (a) Increasing \(b_1\) with \(n_f\) shifts \(d_\text{intersect}\) above the primary particle diameter \(d_p\), while decreasing \(b_1\) shifts it below.
    (b) When evaluating the idealized intersection assuming constant \(b_1\), the lower tail of measured floc diameters extends to the left of \(d_\text{intersect}\). This violates the physical constraint that flocs cannot be smaller than their fundamental building blocks.
    (c) Example where \(d_\text{intersect} > d_p\), emphasizing that \(b_1\) must rise with \(n_f\) to preserve the proper size hierarchy as denser flocs experience greater drag.
    }
    \label{fig:11}
\end{figure}

To illustrate this effect, we plot settling velocity versus floc diameter for two hypothetical fractal groups ($n_f = 1.5$ and $n_f = 2.0$), both with a common primary particle diameter of $10~\mu$m, using Eq.~\eqref{eq:floc_settling} (Fig.~\ref{fig:11}a). For $n_f = 1.5$, we fix $b_1 = 50$, and for $n_f = 2.0$, we vary $b_1$ from 10 to 100. As the shape factor increases with fractal dimension, the intersection diameter $d_{\mathrm{intersect}}$ between the two groups becomes greater than the true primary particle diameter. The opposite holds if $b_1$ decreases with $n_f$. 

We evaluate the calculated geometric intersection diameters ($d_\text{intersect}$) obtained under the standard assumption of a constant shape factor (Fig.~\ref{fig:11}b). While most measured floc diameters lie to the right of $d_\text{intersect}$, the lower tail of the observed floc size distribution inevitably extends to the left (i.e., some observed $d_f < d_\text{intersect}$). Because a floc physically cannot be smaller than its constituent primary particles ($d_f \ge d_p$), the true primary particle diameter must be strictly smaller than all measured floc diameters. Consequently, the assumption that $d_\text{intersect} = d_p$ produces a physical contradiction, requiring instead that the true primary particle size be securely anchored to the left of the constant-$b_1$ intersection ($d_p < d_\text{intersect}$). 

This strict physical constraint implies a necessary relationship between $b_1$ and $n_f$. According to Eq.~\eqref{eq:dpintersection}, the only way to maintain $d_p < d_\text{intersect}$ as $n_f$ increases is if $b_1$ also increases. Such a trend is consistent with theoretical expectations and experimental observations, which suggest that denser, less permeable flocs (i.e., those with higher $n_f$) experience greater hydrodynamic drag, corresponding to higher values of $b_1$ \cite{strom2011explicit}. Importantly, this conclusion does not invalidate the $n_f$ values inferred from the stratified $\log(w_s)$--$\log(d_f)$ slopes, because those slopes are independent of $b_1$; the shape factor enters only the intercept. The sensitivity of $d_\text{intersect}$ to $b_1$ therefore affects only the magnitude of the inferred primary particle diameter, not the fractal dimension itself nor the systematic variation of $d_p$ across experimental conditions.

\subsection{Reconciling Settling Dynamics with Variable Fractal Dimension}
\label{sec:reconciling}

Having established that both $n_f$ (Section~\ref{sec:stratifying}) and $d_p$ (Section~\ref{sec:dp}) vary systematically with experimental conditions, we now test whether this structural variability explains the settling velocity discrepancies identified in Section~\ref{sec:bulk}. We formulate a modified model that retains the same equilibrium $d_f$ prediction but evaluates $w_s \propto d_p^{\,3-n_f}\, d_f^{\,n_f-1}$ with the measured, condition-specific $n_f$ and $d_p$ rather than fixed values. As with the base model, we normalize the prediction by a single least-squares constant and plot it against the measured medians in Fig.~\ref{fig:6}c,d. However, both models share the identical $d_{f,\text{eq}}$ prediction and each contain exactly one free parameter, so any improvement in the modified model arises solely from incorporating the measured structural variability, not from additional fitting degrees of freedom.

As shown in Fig.~\ref{fig:6}c, the modified model successfully reproduces the non-monotonic peak in settling velocity with shear. At low to moderate shear ($u_* = 0.02$ to $0.03$~m\,s$^{-1}$), the rapid densification of flocs (increasing $n_f$) outpaces the reduction in aggregate size, causing $w_s$ to rise despite shrinking $d_f$. Only at high shear ($u_* = 0.04$~m\,s$^{-1}$) does fragmentation dominate over densification, and $w_s$ declines. Similarly, the modified model captures the steep monotonic decrease of $w_s$ with organic matter (Fig.~\ref{fig:6}d): decreasing $n_f$ and $d_p$ with added OM reduce the effective solid fraction of the aggregates, overwhelming the minor increase in $d_f$.

Across all six experiments, the modified model reduces the root-mean-square error in predicted settling velocity by 34\% relative to the base model. This improvement demonstrates that the discrepancies between the constant-$n_f$ framework and our data are not random scatter but arise from systematic structural changes that the base model cannot represent. Properly predicting the settling dynamics of cohesive sediment therefore requires accounting for the dynamic variability of $n_f$ and $d_p$ with environmental forcing.

\section{Discussion}

A central finding of this research is the discrepancy between fractal dimensions derived from aggregated (i.e., pooled) datasets and those obtained from datasets stratified by 2D box-counting dimension. As illustrated by our conceptual model (Fig.~\ref{fig:1}) and supported by experimental results (Figs.~\ref{fig:8} and \ref{fig:9}), aggregating settling data from a floc population with inherent structural variability---meaning a range of true $n_f$ values---can yield a bulk-fitted $n_f$ that overrepresents certain structural subgroups while failing to capture the full diversity within the population. This approach can obscure underlying trends and produce fitted fractal-dimension values that are potentially misleading. The mechanism is analogous to Simpson’s paradox: because distinct structural subpopulations occupy different regions of $w_s$--$d_f$ space, the slope of a pooled regression is not a simple average of within-group slopes and may be dominated by the subpopulation that exerts the greatest leverage in log--log space.

\subsection{Effect of Shear Velocity and Organic Matter}
The settling velocity $w_s$ governs the rate at which flocculated mud deposits from the water column and is therefore the quantity of most direct relevance to sediment transport predictions. For fractal flocs, $w_s \propto d_p^{\,3-n_f}\, d_f^{\,n_f-1}$ (Eq.~\eqref{eq:floc_settling}), which can be rewritten as $w_s \propto \phi \, d_f^2$ by defining the solid volume fraction $\phi = (d_f / d_p)^{n_f - 3}$. In this form, settling is controlled jointly by floc size ($d_f^2$) and effective density ($\phi$), where the latter encapsulates the combined influence of fractal dimension and primary particle size. This decomposition is useful because shear velocity and organic matter affect $d_f$ and $\phi$ in distinct and sometimes opposing ways.

With increasing shear velocity, $d_f$ decreases as stronger turbulence breaks apart aggregates (Fig.~\ref{fig:6}a), while $n_f$ increases (Fig.~\ref{fig:9}a) and the inferred $d_p$ coarsens (Fig.~\ref{fig:10}d), consistent with preferential retention of compact, shear-resistant structures \cite{winterwerp1998simple}. The rising $n_f$ and coarsening $d_p$ increase $\phi$, partially compensating for the shrinking $d_f$. Because these two contributions act in opposing directions on $w_s$, the observed settling response is non-monotonic (Fig.~\ref{fig:6}c). With increasing organic matter, $n_f$ decreases (Fig.~\ref{fig:9}b) and $d_p$ fines (Fig.~\ref{fig:10}h), consistent with organic binding promoting open, porous architectures built from smaller primary particles. Both effects reduce $\phi$, and because $d_f$ remains relatively stable across the OM gradient (Fig.~\ref{fig:6}b), $w_s$ decreases monotonically (Fig.~\ref{fig:6}d).

The equilibrium floc diameter $d_{f,\text{eq}}$, by contrast, is well predicted by the \citeA{nghiem2022mechanistic} framework across all conditions (Fig.~\ref{fig:6}a,b), because $n_f$ cancels in the aggregation--breakup balance and the predicted size depends only on $\eta$ and $\theta$. The settling velocity retains an explicit dependence on $n_f$ and $d_p$, and a constant-$n_f$ model cannot reproduce either the non-monotonic shear response or the organic-matter decline (Fig.~\ref{fig:6}c,d). Incorporating the measured structural variability (Section~\ref{sec:reconciling}) reduces the root-mean-square error by 34\% with no additional free parameters, confirming that these discrepancies reflect systematic structural changes rather than scatter. Although the exact quantitative relationships are likely system-specific, the general trends are broadly consistent with previous studies and underscore that accurate settling predictions require accounting for the variability of $n_f$ and $d_p$ with environmental forcing.

\subsection{2D to 3D Conversion of Fractal Dimension}
The relationship between the 2D box-counting dimension ($n_f^{(\mathrm{2D})}$) and the inferred 3D fractal dimension ($n_f$) is broadly consistent with models developed for synthetic computational aggregates (Fig.~\ref{fig:9}c). The simple Euclidean conversion, $n_f = n_f^{(\mathrm{2D})} + 1$, systematically overestimates $n_f$ across our measured range, whereas our data align more closely with the empirical conversion curves of \citeA{wang2022fractal}, based on computationally generated aggregates. 

However, no universally valid mapping between 2D and 3D fractal dimensions exists for natural flocs. Most synthetic studies, including \citeA{wang2022fractal}, grow aggregates from monodisperse primary particles with fixed size, allowing simple scaling laws to be derived. Natural flocs, by contrast, contain a broad distribution of particle sizes, shapes, and compositions, adding heterogeneity in packing and porosity that is not captured by idealized models. Other approaches \cite{maggi2004method, tang2015reconstructing} depend on their own specific assumptions and likewise cannot account for this variability. These differences introduce scatter and systematic offsets between natural data and synthetic conversions.

The close match between our results and the box-counting conversion of \citeA{wang2022fractal} may therefore reflect structural similarities between our experimental flocs and their synthetic aggregates rather than a universally applicable rule. Nevertheless, the 2D box-counting dimension remains a practical stratification tool: by quantifying within-population variability in $n_f^{(\mathrm{2D})}$ and calibrating the 2D--3D relationship directly from experimental data, our approach provides context-sensitive inference of the hydrodynamically relevant 3D fractal dimension without relying on a fixed conversion law.

\subsection{Constraining Primary Particle Diameter ($d_p$) and Shape Factor ($b_1$)}
The intersection analysis yields $d_\text{intersect}$ distributions that agree in central tendency with the input sediment but are substantially narrower, peaking in the coarse-silt range (Fig.~\ref{fig:10}). This narrowing is expected on theoretical grounds. \citeA{nghiem2024testing} showed that the effective $d_p$ in fractal scaling theory is better represented by a number-weighted moment of the size distribution rather than a volumetric median, because fractal aggregation conserves particle number rather than volume. The inferred $d_\text{intersect}$ is therefore expected to deviate from the bulk median grain size in a manner that depends on the weighting imposed by the aggregation process. Additionally, the finest clays may remain dispersed under the experimental freshwater chemistry, while the coarsest silica grains may be too heavy to be incorporated into flocs---both effects narrow the effective $d_p$ distribution relative to the input. 

The intersection method assumes a constant $b_1$ across $n_f$ classes (Eq.~\eqref{eq:dpintersectionideal}); when $b_1$ varies with floc structure, $d_\text{intersect}$ deviates from the true $d_p$ (Eq.~\eqref{eq:dpintersection}). Our analysis (Fig.~\ref{fig:11}) shows that physically plausible intersections require $b_1$ to increase with $n_f$, consistent with the expectation that denser, less permeable flocs experience greater form drag \cite{strom2011explicit}. However, this sensitivity affects only the absolute magnitude of $d_\text{intersect}$. The $n_f$ values inferred from stratified slopes are independent of $b_1$, which enters only the intercept, and the relative trends in $d_\text{intersect}$ across experimental conditions are preserved because the same $b_1(n_f)$ relationship applies uniformly to all experiments. Independently constraining $b_1(n_f)$---for example through direct drag measurements on size-selected flocs---would allow the intersection estimate to be converted into a true $d_p$. 

Floc permeability, not considered here, may further complicate the inference. Permeability reduces the effective drag coefficient and may itself depend on $d_f$ and $d_p$; if it varies systematically with floc size, assuming it to be constant can bias both the inferred slope and therefore $n_f$ \cite{nghiem2024testing}. Quantifying the joint dependence of $b_1$ and permeability on floc structure remains an important target for future work.

\subsection{Implications}
The 3D fractal dimension of mud flocs, $n_f$, is a key parameter in morphodynamic models of fine-grained sediment transport. Historically, modelers have generally assumed $n_f \approx 2$ \cite{kranenburg1994fractal, son2008flocculation, winterwerp1998simple, winterwerp2004introduction}. These studies inferred the fractal dimension by aggregating settling velocity data from entire floc populations. When we apply this same classical aggregation technique to our own experimental data, we derive a similar fractal dimension of approximately 2.0. This consistency confirms that our experimental flocs are physically comparable to those observed in previous work. However, this apparent agreement masks a significant underlying complexity. 

By stratifying flocs according to their 2D box-counting dimension, our method reveals systematic trends that are invisible to conventional data aggregation. We find that the median $n_f$ increases with shear velocity (from 1.77 to 2.12) and decreases with increasing organic matter content (from 1.84 to 1.63). When the stratified $n_f$ and $d_p$ values are fed back into the settling velocity equation (Section~\ref{sec:reconciling}), the modified model reduces the prediction error by 34\% relative to the constant-$n_f$ base model, providing direct evidence that this structural variability is dynamically significant. In contrast, the aggregated analysis compresses the dynamic range by roughly a factor of five, yielding median values that vary by only $\sim$0.06 across the same experimental conditions compared with $\sim$0.2--0.35 for the stratified approach. The aggregated estimates also regress toward a population-average $n_f \approx 2$, reflecting Simpson’s paradox: fitting a single slope to a structurally mixed population of loose and dense flocs yields an intermediate value that neither represents the dominant subgroup nor captures the underlying environmental trend. The result is that relying on a single, representative fractal dimension near 2.0, as is typical in existing models, both overstates floc compactness under conditions that produce porous flocs and obscures the sensitivity of $n_f$ to shear and organic matter. 

This finding has important implications for sediment transport prediction. Using the Stokes' law--based formulation for floc settling velocity, $w_s \propto d_p^{3-n_f}\, d_f^{n_f-1}$, lowering $n_f$ from 2.0 to 1.7 for a typical 400~\textmu m floc (with 4~\textmu m primary particles, a representative value from the literature) produces a several-fold decrease in settling speed, which directly translates to a several-fold increase in transport distance. In many transport formulations the deposition flux scales as $F_\text{dep} \sim w_s\, C$ (where $C$ is the suspended sediment concentration) \cite{amoudry2008review}, so lowering $n_f$---and thus $w_s$---reduces modeled deposition rates and increases downstream export for the same supply. Flocs that might be predicted to settle out over a given distance would travel several times farther under this revised, more porous structural model. As a result, mud flocs in natural environments may settle more slowly and be transported substantially farther than predicted by existing models that rely on aggregated fractal dimensions. The efficiency of mud retention, at least in the freshwater environments considered here, may be considerably lower than currently assumed, with significant consequences for sediment budgets, contaminant dispersal, and carbon burial in floodplains and deltas.

\section{Conclusion}
This study demonstrates that accounting for heterogeneity within floc populations is essential for accurate sediment transport modeling. Aggregating settling data across flocs with varying fractal dimensions can produce misleading results due to Simpson's paradox: trends present within structural subgroups are lost or distorted when data are pooled. By stratifying flocs based on their 2D box-counting fractal dimension, we resolve structural differences that conventional bulk analysis obscures. 

The stratified approach reveals that $n_f$ increases with shear velocity and decreases with organic matter content, while the inferred $d_p$ co-varies with these same conditions---coarsening under stronger shear and fining under organic binding. The concurrent variation of $n_f$, $d_p$, and $d_f$ introduces competing effects on $w_s \propto d_p^{3-n_f}\, d_f^{n_f-1}$, producing non-monotonic settling responses that constant-$n_f$ models cannot capture. Incorporating the measured structural variability into the settling velocity prediction reduces the root-mean-square error by 34\% relative to the constant-$n_f$ base model, with no additional free parameters. The framework also yields an empirical 2D--3D fractal dimension relationship that matches existing synthetic-aggregate conversions and shows that $b_1$ must increase with $n_f$ to maintain physical consistency. 

These results indicate that the conventional assumption of $n_f \approx 2$ overestimates floc compactness and settling rates. Reducing $n_f$ to the observed range of 1.6--2.1 can decrease predicted $w_s$ by several-fold for typical floc sizes, implying that mud retention in freshwater systems may be considerably lower than current models suggest. We recommend adopting a lower default $n_f$ and, where feasible, representing the fractal dimension as a distribution. While no single value captures the full complexity of natural flocculation, shifting away from the classical range is a necessary step toward more accurate predictions of fine sediment transport, deposition, and downstream sediment budgets.



\section*{Open Research Section}
All code and analysis workflows supporting the results and conclusions of this study are openly available at \url{https://github.com/braydennoh/FlocTrack}. The repository includes scripts for image analysis, fractal dimension calculation, and all processing steps described in this manuscript. For further questions or access to additional datasets, please contact the corresponding author at braydennoh@fas.harvard.edu.

\section*{Inclusion in Global Research Statement}
This research did not involve collaborations with local communities or cross-regional partnerships requiring additional disclosure. All research activities were conducted in accordance with institutional and international ethical standards.

\acknowledgments
BN acknowledges support from Caltech SURF, the Arthur R. Adams SURF Fellowship, and the Saul and Joan Cogen Memorial SURF Fellowship at Caltech. MPL acknowledges support from NASA FINESST Grant 80NSSC20K1645, the NASA Delta-X project funded by the Science Mission Directorate's Earth Science Division through the Earth Venture Suborbital-3 Program NNH17ZDA001N-EVS3, and the National Science Foundation Geomorphology and Land-use Dynamics Grant No. 2136991.

\bibliography{agusample}


\end{document}
